% 本书原创：Green 表示公式中挖孔区域外法向的二维截面示意。
\begin{figure}[htbp]
\centering
\begin{tikzpicture}[x=1cm,y=1cm,>=stealth,line width=.7pt,every node/.style={font=\small}]
\fill[AnalysisBlue!7] (0,0) ellipse (3.3 and 1.65);
\draw[AnalysisBlue] (0,0) ellipse (3.3 and 1.65);
\fill[white] (.55,.05) circle (.65);
\draw[AnalysisBlue] (.55,.05) circle (.65);
\node at (-1.5,.15){\(\Omega_\varepsilon\)};
\node[above] at (-1.8,1.5){\(\partial\Omega\)};
\fill (.55,.05) circle (1.4pt) node[below left]{\(x\)};
\fill (.55,.7) circle (1.4pt) node[above left]{\(y\)};
\draw[densely dashed,gray] (.55,.05)--(1.2,.05) node[midway,below]{\(\varepsilon\)};
\draw[->,AnalysisBlue,thick] (.55,.7)--(.55,.27);
\draw[gray,thin] (.66,.43)--(1.2,.55);
\node[anchor=west] at (1.3,.55){\(\nu_y=-e_r\)};
\draw[->,gray,dashed] (.55,.7)--(.55,1.15) node[above right]{\(e_r\)};
\draw[->,AnalysisBlue,thick] (3.3,0)--(4.05,0) node[right]{\(\nu_\Omega\)};
\end{tikzpicture}
\caption[挖孔区域的内外边界法向]{挖孔区域的内外边界法向。图为 \(d\geq2\) 时的二维截面示意，\(\Omega_\varepsilon=\Omega\setminus\overline{B(x,\varepsilon)}\)，不按比例；内边界的外法向指向孔内，即 \(\nu_y=-(y-x)/\varepsilon\)。由于 \(\Phi_d\) 的径向导数为 \(-1/(\sigma_{d-1}\varepsilon^{d-1})\)，故内球面上 \(\partial_{\nu_y}\Phi_d(x-y)=1/(\sigma_{d-1}\varepsilon^{d-1})\)。边界积分中的 \(-u\,\partial_{\nu_y}\Phi_d\) 恰为负球面平均，积分极限等于 \(-u(x)\)；含 \(\Phi_d\,\partial_\nu u\) 的另一项趋零。}
\label{b2:fig:13-punctured-normal}
\end{figure}
